\documentclass[11pt]{article}

\usepackage[margin=1in]{geometry}
\usepackage{amsmath,amssymb,mathtools}
\usepackage{enumitem}
\usepackage{microtype}
\usepackage[hidelinks]{hyperref}

\newcommand{\E}{\mathbb E}
\newcommand{\indep}{\mathrel{\perp\!\!\!\perp}}
\newcommand{\op}{\mathrm{op}}
\newcommand{\obs}{\mathrm{obs}}

\title{The Graph-HAC Centering Defect and a Feasible Correction}
\author{Exposure-Mapping Project: Current Clarification}
\date{August 3, 2026}

\begin{document}
\maketitle

\section*{Purpose}

This note isolates and repairs the graph-HAC defect on page 21 of
\texttt{park exposure 0730.pdf}. The central problem is not Theorem 2.6: that
theorem and its Appendix A proof are sound. The problem is the later move from
the true parameter \(\theta_0\), where the observed unit score is mean zero, to
a preliminary estimator \(\widehat\theta_n\), where the unit-level score mean
is generally unknown. The proof audit needed to establish this distinction is
relegated to Appendix A.

All expectations are design expectations under \(W\sim D_n\), conditional on
the fixed potential-outcome schedule and baseline covariates. Page references
are to the printed pages of the July 30 draft.

\section{Minimal setup}

Park's Section 2 uses \(i=1,\ldots,n\), the random assignment vector
\(W=(W_1,\ldots,W_n)^\top\), its support \(\mathcal W\), and a fixed assignment
\(w\in\mathcal W\). The three outcome objects that must not be conflated are
\[
 \underbrace{Y_i(w)}_{\text{fixed potential outcome}},\qquad
 \underbrace{\mathsf{Y}_i:=Y_i(W)}_{\text{random realized selector}},\qquad
 \underbrace{Y_i^{\obs}:=Y_i(w^{\obs})}_{\text{observed value}}.
\]
Park writes the middle object simply as \(Y_i=Y_i(W)\). The symbols
\(\mathsf{Y}_i\), \(Y_i^{\obs}\), and \(w^{\obs}\) are introduced only in this
note to make the timeline visible.

Park states the orthogonal moment for a generic outcome transformation
\(\phi\). Because the leading case at issue uses the outcome itself, this note
specializes once and for all to
\[
 \phi(y)=y.
\]
Thus no transformation notation is needed below. This specialization requires
\(\E[\mathsf{Y}_i^2]<\infty\) and gives the same outcome-level score used in the
HAC discussion.

For a candidate \(\theta\), retain Park's design residual notation
\begin{equation}
 R_{i,\theta}\psi(W)
 :=\psi(W)-\E[\psi(W)\mid g_i(W;\theta)],
 \qquad
 g_i(W;\theta):=g(W;X_i,\theta).
 \label{eq:residual}
\end{equation}
In Section 3 Park switches to the triangular array
\(\mathcal U_n=\{1,\ldots,N_n\}\). After the identity specialization, the
relevant scalar unit score is
\begin{equation}
 \Psi_i(\theta):=\mathsf{Y}_iR_{i,\theta}\psi(W).
 \label{eq:score}
\end{equation}
The vector case simply stacks the chosen design functions. Park's theoretical
centering is
\begin{equation}
 Z_{i,n}(\theta)
 :=\Psi_i(\theta)-\E[\Psi_i(\theta)].
 \label{eq:centered-score}
\end{equation}

\section{The HAC defect}

\subsection{What is computable from the design}

For every fixed candidate \(\theta\), the projection
\[
 \E[\psi(W)\mid g_i(W;\theta)]
\]
depends only on the known assignment law, the chosen design function, and the
candidate exposure. Consequently, the residual in \eqref{eq:residual} can be
computed exactly or approximated with placebo assignments. Given the realized
outcome, the uncentered score \(\Psi_i(\theta)\) can therefore be computed for
every candidate \(\theta\).

This does \emph{not} make the separate score mean
\(\E[\Psi_i(\theta)]\) design-computable. For a discrete design,
\begin{align}
 \E[\Psi_i(\theta)]
 &=\E[\mathsf{Y}_iR_{i,\theta}\psi(W)] \\
 &=\sum_{w\in\mathcal W}
   D_n(w)Y_i(w)R_{i,\theta}\psi(w).
 \label{eq:score-mean}
\end{align}
The residual factor is known, but the sum also contains \(Y_i(w)\) at every
assignment in the support. Only \(Y_i(w^{\obs})\) is observed. Equation
\eqref{eq:score-mean} is therefore generally not identified from the known
design and one realized experiment.

\subsection{Why centering is feasible at \(\theta_0\)}

Under Hypothesis 2.5,
\[
 \mathsf{Y}_i=Y_i(W)
 =\widetilde Y_i\{g_i(W;\theta_0)\}
 \quad\text{almost surely}.
\]
Hence \(\mathsf{Y}_i\) is measurable with respect to the true exposure. By
\eqref{eq:residual},
\[
 \E[R_{i,\theta_0}\psi(W)\mid g_i(W;\theta_0)]=0.
\]
The conditional pull-out property then gives
\begin{align}
 \E[\Psi_i(\theta_0)]
 &=\E[\mathsf{Y}_iR_{i,\theta_0}\psi(W)] \\
 &=\E\!\left[
   \mathsf{Y}_i
   \E[R_{i,\theta_0}\psi(W)\mid g_i(W;\theta_0)]
 \right] \\
 &=0.
 \label{eq:truth-zero}
\end{align}
Thus, unit by unit,
\begin{equation}
 Z_{i,n}(\theta_0)=\Psi_i(\theta_0).
 \label{eq:truth-equality}
\end{equation}
This is the valid part of Remark 3.8: at the truth, the feasible uncentered
observed score is already the theoretically centered score.

\subsection{Why the same argument fails at \(\widehat\theta_n\)}

For a generic \(\theta\), the residual remains conditionally mean zero given
\(g_i(W;\theta)\). But \(\mathsf{Y}_i\) need not be measurable with respect to
that candidate exposure. The invalid calculation would be
\[
 \E[\mathsf{Y}_iR_{i,\theta}\psi(W)\mid g_i(W;\theta)]
 \stackrel{\mathrm{not\ valid}}{=}
 \mathsf{Y}_i
 \E[R_{i,\theta}\psi(W)\mid g_i(W;\theta)].
\]
Equivalently,
\begin{equation}
 \E[\Psi_i(\theta)]
 =\E\!\left[
   \operatorname{Cov}\{\mathsf{Y}_i,\psi(W)\mid g_i(W;\theta)\}
 \right],
 \label{eq:off-truth-covariance}
\end{equation}
which need not vanish away from \(\theta_0\).

Consistency of \(\widehat\theta_n\) does not make
\(\E[\Psi_i(\widehat\theta_n)]\) exactly zero, nor does it reveal the missing
potential outcomes in \eqref{eq:score-mean}. It may make the effect of using
\(\widehat\theta_n\) asymptotically negligible, but that is a separate
stochastic-equicontinuity claim that must be assumed or proved.

\subsection{The exact problem in Assumption 3.9}

Park defines \(Z_{i,n}(\theta)\) by \eqref{eq:centered-score}. Assumption 3.9
then describes graph-HAC estimators constructed from the centered variables
and states that, at a preliminary \(\widehat\theta_n\), the design centering is
computed from the known design rather than estimated (source, p.~21). Four
facts must be separated:
\begin{enumerate}[leftmargin=2em]
 \item The \(Y_i\) in the score correctly denotes the random outcome
 \(Y_i(W)\).
 \item The unit score is mean zero at \(\theta_0\), so
 \(Z_{i,n}(\theta_0)=\Psi_i(\theta_0)\).
 \item The projection inside \(R_{i,\theta}\) is design-computable for every
 \(\theta\).
 \item The additional centering \(\E[\Psi_i(\widehat\theta_n)]\) is generally
 neither zero nor design-computable from one realized outcome.
\end{enumerate}
The fourth statement is the defect. The source itself uses the same
missing-potential-outcome logic on pages 24 and 80 when explaining why the
Stage-2 unit means cannot be identified at unrealized exposures.

\subsection{Centered and feasible HAC objects are not identical off truth}

The following temporary notation makes the gap explicit. It is introduced in
this note and is not Park's notation. Define
\[
 m_i(\theta):=\E[\Psi_i(\theta)].
\]
For illustration, represent a normalized graph-HAC construction as
\begin{equation}
 \mathcal H_n(v)
 :=\frac{1}{N_n}
   \sum_{i\in\mathcal U_n}\sum_{j\in A_i}
   k_{ij,n}v_iv_j^\top,
 \label{eq:hac-operator}
\end{equation}
where \(k_{ij,n}\) are known graph or taper weights. Park leaves the particular
HAC construction high level; \(\mathcal H_n\) and \(k_{ij,n}\) are used here
only to display the centering problem.

Call
\[
 \widehat\Omega_n^c(\theta)
 :=\mathcal H_n\bigl(\{Z_{i,n}(\theta)\}_{i\in\mathcal U_n}\bigr),
 \qquad
 \widehat\Omega_n^u(\theta)
 :=\mathcal H_n\bigl(\{\Psi_i(\theta)\}_{i\in\mathcal U_n}\bigr).
\]
The superscripts \(c\) and \(u\), meaning centered and uncentered, are new in
this note. The first object matches the theoretical centered construction but
is infeasible away from the truth. The second uses observed scores and is
feasible. Since \(\Psi_i=Z_{i,n}+m_i\), their exact difference is
\begin{align}
 &\widehat\Omega_n^u(\theta)-\widehat\Omega_n^c(\theta) \\
 &\quad=\frac{1}{N_n}
   \sum_{i\in\mathcal U_n}\sum_{j\in A_i}k_{ij,n}
   \left\{
     Z_{i,n}(\theta)m_j(\theta)^\top
     +m_i(\theta)Z_{j,n}(\theta)^\top
     +m_i(\theta)m_j(\theta)^\top
   \right\}.
 \label{eq:centered-uncentered-gap}
\end{align}
At \(\theta_0\), every \(m_i(\theta_0)=0\), so the difference is exactly zero.
At \(\widehat\theta_n\), it is not exactly zero and its terms cannot generally
be constructed because the \(m_i(\widehat\theta_n)\) are unknown. A valid
repair therefore cannot claim to compute the centered plug-in object. It must
use the feasible uncentered object and control its plug-in error directly.

\subsection{Why holding the observed outcome fixed does not repair centering}

After observing \(w^{\obs}\), holding \(Y_i^{\obs}=Y_i(w^{\obs})\) fixed while
drawing a fresh \(W'\sim D_n\) gives
\begin{align}
 \E_{W'}[Y_i^{\obs}R_{i,\theta}\psi(W')]
 &=Y_i^{\obs}\E_{W'}[R_{i,\theta}\psi(W')] \\
 &=0.
 \label{eq:hybrid-zero}
\end{align}
This equation is correct but answers a different question. The repeated-design
score is
\[
 Y_i(W')R_{i,\theta}\psi(W'),
\]
not \(Y_i(w^{\obs})R_{i,\theta}\psi(W')\). In the former, the selected potential
outcome changes with \(W'\); in the latter, it is artificially held fixed.
Writing the observed number as an unindexed \(Y_i\) cannot turn one experiment
into the other.

A one-unit example makes the distinction immediate. Let
\(W_i\sim\operatorname{Bernoulli}(1/2)\), \(Y_i(0)=0\), \(Y_i(1)=1\), and
\(r(W_i)=W_i-1/2\). Each fixed potential outcome satisfies
\[
 \E[Y_i(0)r(W_i)]=\E[Y_i(1)r(W_i)]=0,
\]
but the random selector \(\mathsf{Y}_i=Y_i(W_i)=W_i\) satisfies
\[
 \E[\mathsf{Y}_ir(W_i)]=\frac14.
\]

\section{Feasible correction}

\subsection{Corrected high-level HAC conditions}

Use the graph-HAC estimator formed from the feasible uncentered scores,
\(\widehat\Omega_n^u(\widehat\theta_n)\). Replace the plug-in claim in
Assumption 3.9 by two distinct requirements:
\begin{align}
 \left\|
   \widehat\Omega_n^u(\theta_0)
   -\frac{\Omega_n(\theta_0)}{N_n}
 \right\|_{\op}
 &=o_p(1),
 \tag{C-HAC1}\label{eq:hac-truth}\\
 \left\|
   \widehat\Omega_n^u(\widehat\theta_n)
   -\widehat\Omega_n^u(\theta_0)
 \right\|_{\op}
 &=o_p(1).
 \tag{C-HAC2}\label{eq:hac-plugin}
\end{align}
Condition C-HAC1 is ordinary truth-point HAC consistency. It is coherent
because \(\Psi_i(\theta_0)=Z_{i,n}(\theta_0)\). Condition C-HAC2 is the missing
plug-in stability requirement. Neither condition asks the researcher to
compute \(\E[\Psi_i(\widehat\theta_n)]\).

Retain Park's population covariance stabilization condition
\begin{equation}
 \left\|
   \frac{\Omega_n(\theta_0)}{N_n}-\Omega
 \right\|_{\op}=o(1).
 \tag{C-HAC3}\label{eq:hac-population}
\end{equation}

\subsection{Why these conditions fix the variance argument}

By the triangle inequality,
\begin{align}
 &\left\|
   \widehat\Omega_n^u(\widehat\theta_n)-\Omega
  \right\|_{\op} \\
 &\quad\le
 \left\|
   \widehat\Omega_n^u(\widehat\theta_n)
   -\widehat\Omega_n^u(\theta_0)
 \right\|_{\op} \\
 &\qquad+
 \left\|
   \widehat\Omega_n^u(\theta_0)
   -\frac{\Omega_n(\theta_0)}{N_n}
 \right\|_{\op}
 +\left\|
   \frac{\Omega_n(\theta_0)}{N_n}-\Omega
 \right\|_{\op} \\
 &=o_p(1).
 \label{eq:hac-triangle}
\end{align}
Thus the feasible plug-in estimator consistently estimates the asymptotic
covariance. If \(\Omega\) is positive definite, its inverse is also consistent
by continuity of matrix inversion. This supplies the covariance step required
for the two-step GMM weight and the overidentification statistic without ever
constructing the infeasible \(Z_{i,n}(\widehat\theta_n)\).

\subsection{What must establish plug-in stability}

Condition C-HAC2 is substantive; consistency of \(\widehat\theta_n\) alone is
not a proof. Let
\[
 \Delta_{i,n}:=\Psi_i(\widehat\theta_n)-\Psi_i(\theta_0).
\]
For the illustrative construction \eqref{eq:hac-operator}, the plug-in
difference is exactly
\begin{align}
 &\widehat\Omega_n^u(\widehat\theta_n)
   -\widehat\Omega_n^u(\theta_0) \\
 &\quad=\frac{1}{N_n}
   \sum_{i\in\mathcal U_n}\sum_{j\in A_i}k_{ij,n}
   \left\{
     \Psi_i(\theta_0)\Delta_{j,n}^\top
     +\Delta_{i,n}\Psi_j(\theta_0)^\top
     +\Delta_{i,n}\Delta_{j,n}^\top
   \right\}.
 \label{eq:plugin-expansion}
\end{align}
Consequently, a primitive proof must control the weighted aggregate on the
right-hand side.

For a smooth exposure map, one route is a local Lipschitz bound for the score,
\[
 \|\Psi_i(\theta)-\Psi_i(\theta_0)\|
 \le L_{i,n}\|\theta-\theta_0\|,
\]
together with moment and graph-weight conditions that keep the normalized
weighted averages of \(L_{i,n}\) and \(\|\Psi_i(\theta_0)\|\) bounded. The
required rate is the rate that makes all three terms in
\eqref{eq:plugin-expansion} \(o_p(1)\); it must account for growth of the
affinity sets and cannot be replaced by bare consistency.

For the nonsmooth ring map, a Lipschitz argument is unavailable. Here
\(\Delta_{i,n}\) changes only for units whose distance thresholds lie between
\(\theta_0\) and \(\widehat\theta_n\). A sufficient route is a
degree-weighted boundary-mass condition: the normalized HAC weight attached to
pairs affected by that shrinking interval must be \(o_p(1)\). This is the HAC
analogue of the no-mass-in-shrinking-distance-bands condition used for the
ULLN, but it must be verified for the pairwise graph-HAC aggregate.

\subsection{Suggested replacement for the page 21 claim}

A concise replacement for the end of Remark 3.8 and Assumption 3.9 is:
\begin{quote}
At the true exposure parameter, Corollary 2.7 implies that the uncentered unit
score is mean zero, so it coincides with the centered score. Let
\(\widehat\Omega_n^u(\theta)\) be a graph-HAC estimator constructed from the
observed uncentered scores \(\{\Psi_i(\theta)\}_{i\le N_n}\) and the affinity
sets \(\{A_i\}_{i\le N_n}\). Assume C-HAC1 and C-HAC2. Together with
\(\Omega_n(\theta_0)/N_n\to\Omega\), these conditions imply
\(\widehat\Omega_n^u(\widehat\theta_n)\to_p\Omega\).
\end{quote}
The sentence claiming that plug-in score centering is computed from the known
design should be deleted. Only the conditional projection inside
\(R_{i,\theta}\) has that property for arbitrary \(\theta\).

\section*{Conclusion}

The HAC defect is localized and repairable. Theorem 2.6 correctly makes the
score mean zero at \(\theta_0\), but it does not make the score mean known at
\(\widehat\theta_n\). The feasible correction is to form graph-HAC from the
uncentered observed scores, establish HAC consistency at the truth, and add a
separate plug-in stability condition. The triangle argument
\eqref{eq:hac-triangle} then restores covariance consistency.

\appendix

\section{Supporting proof audit}

\subsection{Theorem 2.6}

This appendix records the proof only to verify the outcome object used in the
HAC argument. Assumption 2.4 fixes the schedule
\(\{Y_i(\cdot),X_i\}_{i=1}^n\) relative to the design. Hypothesis 2.5 states
that there are \(\theta_0\) and fixed measurable maps \(\widetilde Y_i\) such
that, for every fixed \(w\in\mathcal W\),
\[
 Y_i(w)=\widetilde Y_i\{g(w;X_i,\theta_0)\}.
\]
Using Park's Appendix A shorthand
\(G_i(\theta_0):=g(W;X_i,\theta_0)\), substitute the random assignment into
this fixed-assignment identity:
\begin{equation}
 \mathsf{Y}_i=Y_i(W)
 =\widetilde Y_i\{G_i(\theta_0)\}
 \quad\text{almost surely}.
 \label{eq:appendix-measurability}
\end{equation}
Because \(\widetilde Y_i\) is fixed under the design, \(\mathsf{Y}_i\) is
measurable with respect to \(\sigma\{G_i(\theta_0)\}\). For bounded measurable
test functions \(f\) and \(h\), the conditional pull-out property gives
\begin{align}
 &\E[f(\mathsf{Y}_i)h(W)\mid G_i(\theta_0)] \\
 &\quad=f\{\widetilde Y_i(G_i(\theta_0))\}
       \E[h(W)\mid G_i(\theta_0)] \\
 &\quad=\E[f(\mathsf{Y}_i)\mid G_i(\theta_0)]
       \E[h(W)\mid G_i(\theta_0)].
\end{align}
This factorization is exactly
\[
 \mathsf{Y}_i\indep W\mid G_i(\theta_0).
\]
The conditional-independence equation therefore comes from the exposure-map
restriction and measurability, with Assumption 2.4 fixing the response maps
under the design. Random assignment alone would not imply it.

\subsection{Corollary 2.7 under the identity specialization}

By definition,
\[
 \E[R_{i,\theta_0}\psi(W)\mid G_i(\theta_0)]=0.
\]
Equation \eqref{eq:appendix-measurability} makes \(\mathsf{Y}_i\) measurable
with respect to the conditioning variable. Under square integrability,
\begin{align}
 \E[\mathsf{Y}_iR_{i,\theta_0}\psi(W)]
 &=\E\!\left[
   \E[\mathsf{Y}_iR_{i,\theta_0}\psi(W)\mid G_i(\theta_0)]
 \right] \\
 &=\E\!\left[
   \mathsf{Y}_i
   \E[R_{i,\theta_0}\psi(W)\mid G_i(\theta_0)]
 \right] \\
 &=0.
\end{align}
This is the equality used in \eqref{eq:truth-zero}. In Park's notation it is
\(\E[Y_iR_{i,\theta_0}\psi(W)]=0\), with the unindexed \(Y_i\) denoting the
random selector \(Y_i(W)\), not a fixed \(Y_i(w)\).

\end{document}
